\input{preamble.tex}

\begin{document}

%%%%%%%%%%%%%%%%%%%%%%%%%%%%%%%%%%%%%%%%%%%%%%%%%%%%%%%%%%%%
    
\input{frontpage.tex}

%%%%%%%%%%%%%%%%%%%%%%%%%%%%%%%%%%%%%%%%%%%%%%%%%%%%%%%%%%%%

\tableofcontents
\listoffigures
%\listoftables

%%%%%%%%%%%%%%%%%%%%%%%%%%%%%%%%%%%%%%%%%%%%%%%%%%%%%%%%%%%%
% 1. Presentació
%%%%%%%%%%%%%%%%%%%%%%%%%%%%%%%%%%%%%%%%%%%%%%%%%%%%%%%%%%%%
\chapter{Presentació}
Una breu presentació del que fareu en aquest diari de treball.

\section{Conceptes bàsics}

\subsection{Confort ambiental}

\subsection{Confort tèrmic}

\subsection{Illa de calor}


\section{Equip i primera proposta de treball}
Aquí podríeu parlar de cóm heu conformat parella i per què (proximitat geogràfica, confiança en la capacitat de treball, habilitats complementaries, amistat, etc). També podríeu reflexionar sobre una zona d'estudi i si hi ha activitats a la zona que puguen estar relacionades amb el confort tèrmic.


\newpage

%%%%%%%%%%%%%%%%%%%%%%%%%%%%%%%%%%%%%%%%%%%%%%%%%%%%%%%%%%%%
% 2. Literatura
%%%%%%%%%%%%%%%%%%%%%%%%%%%%%%%%%%%%%%%%%%%%%%%%%%%%%%%%%%%%

\chapter{Revisió de literatura científica}
Aquí cap una breu reflexió sobre el que preteneu amb la consulta bibliogràfica.

\section{Eines utilitzades}
Mendeley, Overleaf, etc

\begin{figure}
    \centering
    \includegraphics[width=\textwidth]{images/template/image.png}
    \caption{Les imatges aniràn colocades allà on \LaTeX{} considere més apropiat però també podeu forçar-ho. De moment com hi ha poc text poc resultar una mica estrany. En aquest cas està definit que la imatge ocupe tot l'ample del text.}
    \label{fig:mendeleydesktop}
\end{figure}{}

\section{Bases de dades consultades}
WOS, Google scholar, etc

\section{Síntesi dels treballs revisats}
Temàtiques, problemes estudiats, metodologies que apareixen als documents consultats.Haureu de citar les vostres fonts amb el comandament corresponent. Per exemple, \cite{PerezCueva2006} per citar explícitament o \citep{PerezCueva2006}, si voleu parèntesis.


%%%%%%%%%%%%%%%%%%%%%%%%%%%%%%%%%%%%%%%%%%%%%%%%%%%%%%%%%%%%
% 3. Dades
%%%%%%%%%%%%%%%%%%%%%%%%%%%%%%%%%%%%%%%%%%%%%%%%%%%%%%%%%%%%
\chapter{Descàrrega i preparació de dades}
Explicar la tasca de descàrrega, quina part de les dades heu preparat, temps que heu tardat en fer el procés d'obtenir un full de càlcul, etc

%%%%%%%%%%%%%%%%%%%%%%%%%%%%%%%%%%%%%%%%%%%%%%%%%%%%%%%%%%%%
% Conclusions
%%%%%%%%%%%%%%%%%%%%%%%%%%%%%%%%%%%%%%%%%%%%%%%%%%%%%%%%%%%%

\chapter{Conclusions}
A partir de la literatura consultada, expliqueu breument el que creieu que deuríeu aprendre per a ver un anàlisi de confort ambiental (tècniques), les dades que caldria recollir, dimensions de la zona d'estudi, etc

\section{Comentaris}
Amb tot, tracteu de centrar el tema que podríeu treballar, una zona d'estudi, l'interès i els destinataris dels vostres resultats. Recordeu que aneu a caracteritzar el confort tèrmic d'una zona i que podríeu aplicar-ho a un col·lectiu o activitat concreta.

\section{Opinió personal}


%%%%%%%%%%%%%%%%%%%%%%%%%%%%%%%%%%%%%%%%%%%%%%%%%%%%%%%%%%%%

\appendix
% Els annexes complementen la informació presentada en el cos principal de l’informe.

% Annex A:
\chapter{Mini-tutorial de Mendeley}
Alguna captura de pantalla o passos que no voleu oblidar.

\begin{figure}
    \centering
    \includegraphics[width=10cm]{images/template/image.png}
    \caption{Podria ser una captura de pantalla. Podeu definir les dimensions exactes de la imatge com es mostra en aquest exemple.}
    \label{fig:mendeleydesktop2}
\end{figure}{}


%%%%%%%%%%%%%%%%%%%%%%%%%%%%%%%%%%%%%%%%%%%%%%%%%%%%%%%%%%%%

% Annex B:
\chapter{Xulla de \LaTeX{}}
Estructura d'un document, els entorns i comandos més habituals.

\begin{itemize}
    \item \verb|\chapter{}| Per marcar el començament d'un nou capítol.
    \item \verb|\section{}| Per marcar el començament d'una nova secció dins un capítol.
    \item \verb|\label{}| Per anotar qualsevol element o part del text i després fer-ne referència amb \verb|\ref{}|.
    \item Haureu de buscar alguna ``xulla'' de \LaTeX{} o fer algun tutorial. Cal treballar elements sobre com inserir imatges, taules, llistes numerades a no numerades, etc

\end{itemize}{}

%%%%%%%%%%%%%%%%%%%%%%%%%%%%%%%%%%%%%%%%%%%%%%%%%%%%%%%%%%%%

\addcontentsline{toc}{chapter}{Bibliografia}
\nocite{*} % Afegeix tots els documents de la bibliografia
\bibliographystyle{apacite}
\bibliography{bibliography/projecte3.bib}


\end{document}
